\paragraph{Locate each declaration in the source that was supplied.}
The matched text extension now expresses these same bounded functions. For
example, the ordinary collection fraction is written as follows:
\begin{lstlisting}[basicstyle=\ttfamily\scriptsize,breaklines=true]
function complement;
parameter discount '{"context":null,"dimensions":{},"kind":"number"}';
returns '{"context":null,"dimensions":{},"kind":"number"}';
param_ref d discount;
literal one '{"context":null,"dimensions":{},"kind":"number"}' 1/1;
node r sub one d;
result r;
end_function;
\end{lstlisting}
The quoted JSON type denotes a number with no dimensions and no company
context. The parameter reference binds the local name \code{d} to the declared
discount; \code{r} subtracts it from one. The result declaration identifies
the returned node, and \code{end\_function} closes the function. In the
surrounding program, \code{call regular complement discount;} supplies the
main program's discount node and names the resulting call \code{regular}.
Input declarations and the other functions remain necessary for the full
collection policy. The displayed wrapping does not introduce extra source
lines: each declaration occupies one physical line and ends in a semicolon.

Function parameters precede the return type and body; a result ends the body.
Blocks cannot nest, but a call can name a function declared later. Main inputs,
nodes and outputs use the earlier expression notation, within a
\code{CompanyFunctions/1} program ending in \code{end\_program}.
The parser rejects duplicate names and JSON keys, incomplete blocks, unsupported
declarations and semantic errors reported by the pinned typed compiler.
Quoted command-looking characters are text: parsing performs no shell
expansion, arithmetic evaluation or event generation. A validly typed division
by zero can therefore compile and later fail when its value is requested.

Each accepted declaration records its path within the typed request, physical
line and exact original UTF-8 byte interval. Byte offsets are zero-based and
end-exclusive; line and Unicode-codepoint columns are one-based, with the ending
column exclusive. The slice includes the terminal semicolon and excludes
surrounding spaces, tabs and the line ending. Its hash and the complete source
hash bind these locations to the supplied bytes. This distinction matters when
non-ASCII names use several bytes per character, or a source uses CRLF rather
than LF. Both line-ending forms are accepted without Unicode normalization.

Canonical rendering uses LF and deterministic quoting. It reproduces the same
typed request but can have a different source hash and different locations.
Consequently, original locations are never assigned to re-rendered text.
The generated expression's function and node origins can be joined to these
declarations; token-level debugging and runtime exception locations remain
separate work. Source locations also supply no evidence of when a policy was
written, published or used.

The text frontend bounds source size to 524,288 bytes, each physical line to
16,384 bytes and accepted declarations to 12,000 before the typed compiler's
own limits apply. These are scoped engineering limits. Eighteen focused test
methods, 23 typed-function regression methods, 79 expression checks and 285
source/financial comparisons cover six policy cases, preserving complete
three-year results where available and unresolved or rejected paths elsewhere.
The exact examples, declaration slices and financial comparisons are retained
in \path{data/accounting-master/company-function-text-20261025/}. The checked
extension still does not supply dates, polymorphism, module dispatch or
endogenous economic decisions.

